\documentclass[11pt]{article}

\usepackage[T1]{fontenc}
\usepackage[utf8]{inputenc}
\usepackage{lmodern}
\usepackage[margin=1in]{geometry}
\usepackage{microtype}
\usepackage{enumitem}
\usepackage{tabularx}
\usepackage{pdflscape}
\usepackage[round,authoryear]{natbib}
\usepackage[hidelinks]{hyperref}
\renewcommand{\abstractname}{Resumo}
\renewcommand{\refname}{Referências}

\title{Axodus: Uma Arquitetura Conceitual em Estágio de Protótipo para Rastreabilidade e Responsabilização em Instituições Digitais Modulares Federadas}
\author{Mauricio Z. Filho\\
\small Pesquisador independente, Axodus\\
\small Correspondência: mzfshark@gmail.com\\
\small Discussão do projeto: Telegram @thinkincoin\\
\small Perfil do autor: @mzfshark}
\date{Rascunho v0.6}

\begin{document}
\hypersetup{
  pdftitle={Artigo de Arquitetura Axodus v0.6 - Tradução Controlada pt-BR},
  pdfauthor={Mauricio Z. Filho},
  pdfsubject={Tradução controlada privada para revisão acadêmica; distribuição não autorizada},
  pdfkeywords={arquitetura conceitual, revisão acadêmica controlada, tradução português-BR}
}
\maketitle

\begin{center}
\fbox{\parbox{0.88\linewidth}{\centering
\textbf{Status:} Rascunho acadêmico interno v0.6 preparado como cópia privada para revisão humana controlada; a entrega a revisores não está autorizada, e a circulação externa e a publicação permanecem bloqueadas.\\
\textbf{Tipo de documento:} Arquitetura conceitual em estágio de protótipo e artigo com agenda de pesquisa focalizada.\\
\textbf{Licença:} CC BY 4.0.}}
\end{center}

\begin{center}
\fbox{\parbox{0.92\linewidth}{\centering
\textbf{CÓPIA DE REVISÃO CONTROLADA - TRADUÇÃO PORTUGUÊS-BR}\\
Este documento é uma tradução controlada para o português brasileiro do
manuscrito autoritativo em inglês, preparada para revisão acadêmica privada.
Em caso de divergência interpretativa, prevalece a versão inglesa verificada.
ID do pacote: \texttt{AXODUS-V0.6-CHR-PTBR-20260712}. Não autorizada para
circulação pública, republicação, encaminhamento, publicação, submissão,
apresentação a parceiros, captação de recursos, uso por investidores ou uso em
treinamento. A entrega a revisores exige autorização nominal separada. A
revisão controlada não implica endosso, validação, aceitação ou conclusão de
revisão por pares.}}
\end{center}

\begin{abstract}
Instituições digitais assistidas por humanos e IA tornam-se difíceis de
governar quando restrições constitucionais, domínios delimitados de governança
local, capacidades modulares e obrigações de revisão precisam ser coordenados
no mesmo ambiente decisório. Literaturas adjacentes sobre sistemas
sociotécnicos, governança descentralizada, proveniência, algoritmos
responsabilizáveis, contestabilidade e interação humano--IA oferecem
componentes com escopos distintos para examinar esse problema. Este artigo
apresenta o Axodus como uma integração proposta desses componentes: uma
arquitetura conceitual em estágio de protótipo para esse contexto. Sua
contribuição central é um modelo de seis camadas orientado à governança que
relaciona um centro constitucional, domínios delimitados de governança local,
capacidades institucionais modulares, evidência e proveniência, assistência de
IA delimitada e controles de risco ou contestabilidade. As seis camadas são
apresentadas como uma decomposição conceitual defensável, não como estrutura
necessária, mínima, unicamente correta ou universalmente suficiente.
Rastreabilidade e reconstrutibilidade podem apoiar responsabilização
(\emph{accountability}) e contestabilidade, mas não estabelecem nenhuma dessas
propriedades de forma independente. O Axodus é descrito somente como um modelo
conceitual em estágio de protótipo; o artigo não relata implantação, validação
empírica, uso institucional observado ou prontidão operacional.
\end{abstract}

\noindent\textbf{Palavras-chave:} sistemas sociotécnicos; instituições
digitais federadas; coordenação humano--IA; rastreabilidade;
responsabilização; governança; contestabilidade

\section{Introdução}

Instituições digitais combinam cada vez mais autoridade distribuída,
capacidades modulares, registros públicos e assistência automatizada. A
governança torna-se mais difícil quando esses elementos precisam permanecer
revisáveis entre um centro constitucional e domínios delimitados de governança
local, em vez de permanecerem dentro de uma única cadeia de autoridade. Nesse
contexto, o desafio institucional não é apenas saber se uma decisão pode ser
documentada, mas se revisores posteriores conseguem reconstruir qual
autoridade se aplicava, qual capacidade delimitada foi invocada, quais
evidências foram consideradas, como a assistência de IA moldou o registro e
como uma correção ou contestação poderia prosseguir. A pesquisa sobre sistemas
sociotécnicos trata essas questões como simultaneamente organizacionais e
técnicas, e não como redutíveis a um único domínio
\citep{baxter2011sociotechnical}.

Literaturas relevantes abordam partes desse problema com escopos distintos e
parcialmente sobrepostos. A pesquisa sobre governança descentralizada mostra
que regras autoexecutáveis não eliminam interpretação, emenda ou coordenação
social \citep{hassan2021dao,hsieh2018bitcoin}. Diretrizes de interação
humano--IA enfatizam revisão, correção e definição de expectativas, em vez de
uma transferência indiferenciada de autoridade
\citep{amershi2019guidelines}. Padrões de proveniência e credenciais ajudam a
distinguir a linhagem da evidência e a verificação por máquina do julgamento
institucional \citep{w3c2013provdm,w3c2025vcdm}. Este artigo não alega uma
demonstração exaustiva de que não existem arquiteturas comparáveis. Ele
desenvolve uma síntese integrativa delimitada para um problema institucional
específico.

Este artigo apresenta o Axodus como uma arquitetura conceitual em estágio de
protótipo para instituições digitais assistidas por humanos e IA sob
restrições delimitadas de governança. O projeto se origina de um esforço
independente de pesquisa e design desenvolvido ao longo de vários anos; essa
afirmação registra a proveniência contextual, e não evidência científica. A
contribuição permanece conceitual: não relata uma plataforma em uso
institucional real, um ambiente institucional em funcionamento, resultados de
avaliação ou controles testados independentemente.

A contribuição consiste em raciocínio de design no nível da arquitetura, e
não em uma alegação de implementação. O artigo propõe um modelo em camadas
para organizar um centro constitucional, domínios delimitados de governança
local, capacidades institucionais modulares, pacotes de evidências,
assistência de IA atribuível e caminhos de revisão em uma estrutura decisória
inspecionável. A estrutura de seis camadas é uma decomposição conceitual
defensável entre alternativas possíveis. A Prontidão de Participação permanece
uma preocupação de apoio delimitada, enquanto as distinções entre
documentação, desenvolvimento, produção e autoridade servem apenas como
limite de escopo.

Implementação, implantação, uso institucional e validação empírica estão fora
do escopo do artigo. O manuscrito também exclui design de tokens, operações
financeiras, mecânicas de distribuição educacional, garantias de segurança e
conclusões jurídicas.

\section{Fundamentação e Trabalhos Relacionados}

\subsection{Sistemas sociotécnicos e instituições digitais federadas}

Uma perspectiva sociotécnica pergunta como papéis, incentivos, práticas de
trabalho, informação e componentes técnicos moldam resultados em conjunto.
Baxter e Sommerville defendem a integração de preocupações organizacionais à
engenharia de sistemas, em vez de tratá-las como um problema tardio de adoção
\citep{baxter2011sociotechnical}. O Axodus adota essa perspectiva ao modelar
governança, responsabilização e revisão como elementos arquiteturais, e não
como políticas externas anexadas após a implementação.

O termo \emph{instituição digital modular federada} é usado aqui para um
sistema organizado no qual uma camada constitucional central define
restrições compartilhadas, enquanto domínios delimitados de governança local
podem configurar capacidades institucionais modulares. Esse uso é conceitual.
Ele não implica uma federação implantada, status federal jurídico ou uma
plataforma multi-inquilino em operação.

\subsection{Governança descentralizada, acoplamento modular e revisão}

Hassan e De Filippi descrevem uma DAO como um sistema baseado em blockchain no
qual pessoas se coordenam por meio de governança descentralizada e regras
autoexecutáveis \citep{hassan2021dao}. Hsieh et al.\ distinguem consenso de
máquina do consenso social necessário para alterar protocolos e arranjos
organizacionais \citep{hsieh2018bitcoin}. Essas distinções recomendam cautela
contra equiparar distribuição técnica a governança responsabilizável.

Mercados financeiros baseados em contratos inteligentes oferecem um exemplo
específico de domínio no qual a composabilidade pode propagar dependências
entre protocolos conectados \citep{schar2021defi}. A pesquisa de segurança
baseada em teoria de sistemas oferece a afirmação mais ampla necessária aqui:
falhas em sistemas sociotécnicos complexos podem surgir de interações,
incompatibilidades de controle e acoplamentos mal governados, e não apenas de
defeitos isolados em componentes \citep{leveson2012engineering}. Por isso, este
artigo trata esse domínio somente como um exemplo motivador de propagação de
dependências, não como seu centro institucional.

\subsection{Supervisão humano--IA e apoio responsabilizável à decisão}

Diretrizes de interação humano--IA enfatizam a comunicação das capacidades do
sistema, o apoio à correção e o aprendizado a partir do comportamento ao longo
do tempo \citep{amershi2019guidelines}. No nível organizacional, o NIST AI Risk
Management Framework trata governança, mapeamento, medição e gestão como
funções de risco conectadas \citep{tabassi2023airmf}. Essas fontes motivam uma
arquitetura na qual saídas de IA permanecem como insumos atribuíveis e
revisáveis para decisões humanas, e não como decisões soberanas em si mesmas.

A IA Constitucional usa princípios escritos para orientar o comportamento do
modelo por meio de aprendizado supervisionado e feedback de IA
\citep{bai2022constitutional}. O Axodus usa \emph{constitucional} de modo
diferente: o termo se refere a regras institucionais sobre autoridade,
proibições, emendas, revisão e contestação. O modelo proposto não alega treinar
ou alinhar um modelo de IA por meio da IA Constitucional.

\subsection{Proveniência, credenciais e julgamento institucional}

O Modelo de Dados PROV do W3C descreve proveniência por meio de entidades,
atividades, agentes, derivações, atribuição e relações de responsabilidade
\citep{w3c2013provdm}. Esse vocabulário é útil para distinguir artefatos de
origem, transformações, assistência automatizada e papéis responsáveis em um
registro decisório. Ele apoia a reconstrução de como uma saída foi produzida,
mas não estabelece, por si só, que a decisão foi legítima, contestável ou
responsabilizável.

O Modelo de Dados de Credenciais Verificáveis do W3C separa a verificação
criptográfica de uma credencial da validação de que suas afirmações são
adequadas ao propósito de um verificador \citep{w3c2025vcdm}. Isso importa
porque artefatos de elegibilidade verificáveis por máquina não estabeleceriam,
por si só, avaliação justa, autorização legítima ou limites de participação
responsabilizáveis.

O trabalho sobre algoritmos responsabilizáveis também argumenta que processos
decisórios automatizados precisam ser avaliáveis em relação a padrões
jurídicos ou de política externos, e que transparência isolada é insuficiente
\citep{kroll2017accountable}. A proveniência é, portanto, um insumo para a
responsabilização, não seu substituto.

\subsection{Contestabilidade e governança revisável}

A literatura sobre contestabilidade por design identifica salvaguardas,
revisão e intervenção humanas e escrutínio de terceiros como características
de design para contestação e correção \citep{alfrink2023contestable}. Esse
enquadramento é diretamente relevante para a governança modular federada,
porque um registro rastreável permanece insuficiente se uma pessoa afetada por
uma decisão não puder identificar a regra aplicável, contestar evidências ou
alcançar um revisor adequadamente independente.

\subsection{Síntese conceitual e derivação da arquitetura}

A arquitetura é derivada por meio de uma síntese conceitual narrativa e não
sistemática, e não por uma revisão sistemática, prova formal ou método de
design empírico. A síntese trata as vertentes da literatura nas Seções
2.1--2.5 como fundamentos conceituais, as cinco tensões na Seção 3 como
problemas de design e os requisitos resultantes como razões para separar
responsabilidades entre seis camadas. As fronteiras entre camadas são traçadas
onde autoridade, interpretação local, capacidade delimitada, reconstrução de
registros, análise consultiva e condições de contestação exigem
responsabilidades diferentes ou poderes expressamente excluídos.

Esse procedimento não estabelece que seis camadas sejam necessárias, mínimas,
unicamente corretas ou universalmente suficientes. Ele apresenta uma
decomposição conceitual defensável cujas fronteiras podem ser inspecionadas,
comparadas com alternativas e revisadas ou rejeitadas. A matriz de derivação
torna explícita a cadeia entre fundamentos e construtos de avaliação; ela não
implica que a literatura citada valide o Axodus.

\begin{landscape}
\begin{center}
\scriptsize
\begin{tabularx}{\linewidth}{|X|X|X|X|X|X|}
\hline
\textbf{Fundamento} & \textbf{Tensão} & \textbf{Requisito} &
\textbf{Camada(s)} & \textbf{Propriedade apoiada} &
\textbf{Construto de avaliação} \\
\hline
Sistemas sociotécnicos e governança institucional & Restrições centrais versus
autonomia local & Separar restrições de ordem superior da interpretação local
delimitada & Governança constitucional; domínio delimitado de governança local
& Condições inspecionáveis de autoridade e prestação de contas & Identificação
de autoridade; detecção de conflitos e de escopo \\
\hline
Modularidade e acoplamento de sistemas & Capacidades modulares versus
governabilidade & Declarar fronteiras de capacidade, dependências e restrições
de design compartilhadas & Módulo funcional de serviço; governança
constitucional & Coerência entre camadas e responsabilidade delimitada &
Completude de dependências; detecção de transições sem suporte \\
\hline
Proveniência e apoio responsabilizável à decisão & Documentação versus
realidade institucional & Atribuir evidências, transformações, papéis e
disposições sem equiparar registros a legitimidade & Evidência e proveniência &
Rastreabilidade, reconstrutibilidade e revisabilidade & Completude da
reconstrução; visibilidade da proveniência; detecção de evidência ausente \\
\hline
Interação humano--IA e governança de risco de IA & Automação versus disposição
responsabilizável & Manter as saídas de IA atribuíveis e consultivas, com
disposição humana explícita & Assistência de IA; evidência e proveniência &
Confiança delimitada e responsabilidade atribuível & Detecção de erros;
comportamento de anulação; qualidade da justificativa \\
\hline
Contestabilidade por design & Disputas de autoridade ou evidência versus
revisão efetiva & Preservar gatilhos de contestação, condições de independência
e caminhos de revisão alcançáveis & Limite de risco e contestabilidade;
governança constitucional & Condições de apoio à contestabilidade e à correção
& Acessibilidade do caminho de contestação; discordância entre revisores;
classificação da disposição \\
\hline
Fronteira de participação e prontidão & Participação versus prontidão
delimitada & Tratar prontidão como insumo contestável e sujeito a adaptações,
não como classificação ou autoridade universal & Domínio delimitado de
governança local; limite de risco e contestabilidade & Condições revisáveis de
participação & Análise de risco de exclusão; visibilidade de adaptações e
contestação \\
\hline
\end{tabularx}
\end{center}
\end{landscape}

\section{Enunciado do Problema}

A arquitetura aborda um problema de design, e não um efeito empírico
estabelecido:

\begin{quote}
Como uma instituição digital modular federada e constitucionalmente restrita
pode apoiar rastreabilidade, reconstrutibilidade e as condições institucionais
para responsabilização e contestabilidade quando domínios delimitados de
governança local configuram capacidades institucionais modulares por meio de
processos decisórios assistidos por humanos e IA?
\end{quote}

Cinco tensões estruturam o problema.

\begin{enumerate}[leftmargin=*]
\item \textbf{Restrições centrais versus autonomia local.} Uma camada
constitucional pode preservar coerência, enquanto domínios delimitados de
governança local ainda exigem discricionariedade significativa sobre decisões
e configurações locais.
\item \textbf{Capacidades modulares versus governabilidade.} A separação de
serviços pode limitar o acoplamento direto, enquanto a seleção modular pode
criar cadeias ocultas de dependência, autoridade ambígua e caminhos
fragmentados de revisão.
\item \textbf{Automação versus disposição responsabilizável.} A assistência
de IA pode reduzir a carga analítica, mas a responsabilidade pode se tornar
incerta quando suas saídas influenciam autorização, escalonamento ou revisão.
\item \textbf{Participação versus prontidão delimitada.} A participação aberta
pode apoiar legitimidade, enquanto ações de governança consequentes podem
exigir preparação revisável, adaptações e contestação.
\item \textbf{Documentação versus realidade institucional.} Uma especificação
detalhada pode melhorar a inspeção e, ao mesmo tempo, criar uma impressão
enganosa de completude de implementação ou autoridade legítima.
\end{enumerate}

O Axodus é proposto como uma decomposição conceitual defensável para tornar
essas tensões inspecionáveis e revisáveis. Ele não é apresentado como solução
completa, ótima, unicamente correta ou testada empiricamente.

\section{Arquitetura Conceitual}

O modelo compreende seis camadas lógicas selecionadas pela derivação da Seção
2.6. Elas são superfícies conceituais de design para alocar autoridade,
estruturar registros e preservar revisabilidade, não uma topologia de
implementação. Sua separação expressa responsabilidades distintas e poderes
excluídos; não afirma que toda instituição exija essa decomposição.

A informação e o controle percorrem a arquitetura de modos diferentes. O
controle começa com restrições constitucionais que definem quais decisões
locais, seleções de capacidades e caminhos de escalonamento estão no escopo.
Pacotes de evidências, registros de proveniência, análise consultiva,
justificativa e registros de revisão tornam essas decisões reconstruíveis para
revisão posterior. Limites de risco restringem transições de estado
permitidas, enquanto caminhos de contestabilidade preservam condições de
contestação, correção, escalonamento e revisão.

\noindent\textbf{Cenário sintético ilustrativo - revisão de exceção de
governança delimitada.} Um domínio simulado de governança local delimitada
solicita uma exceção temporária que afeta uma capacidade institucional modular.
A solicitação inclui a regra constitucional aplicável, a base de autoridade
local, conflitos declarados e um pacote controlado de evidências. Uma revisão
simulada por IA sinaliza um conflito, mas omite uma dependência entre a
configuração solicitada e um requisito de revisão. O revisor humano designado
retém a disposição, registra a omissão e escalona a questão porque um limite de
risco foi acionado. Um papel de revisão independente pode então reconstruir
qual restrição constitucional se aplicava, qual escolha local foi proposta, o
que a revisão simulada por IA alterou e por que ocorreu o escalonamento. Esse
cenário é uma ilustração analítica, não evidência de uma implementação ou de um
processo institucional do Axodus.

\begin{enumerate}[leftmargin=*]
\item a governança constitucional define papéis, direitos de decisão,
proibições, regras de emenda e condições de contestação aplicáveis;
\item domínios delimitados de governança local interpretam esse escopo para um
contexto decisório local e podem propor uma configuração ou exceção dentro dos
limites constitucionais;
\item capacidades institucionais modulares expressam o que está sendo
selecionado, revisado ou alterado;
\item a camada de evidência e proveniência conecta materiais de origem,
declarações de autoridade, conflitos, transformações e papéis responsáveis;
\item a assistência humano--IA pode analisar esse material, mas somente como
insumo consultivo atribuível à disposição humana; e
\item controles de limite de risco e contestabilidade mantêm disponíveis
condições de contestação, correção, escalonamento e revisão.
\end{enumerate}

A tabela de interfaces entre camadas especifica responsabilidades e
transferências conceituais. Seus insumos e saídas são categorias de registros,
não APIs, componentes de software ou alegações sobre infraestrutura
implementada.

\begin{landscape}
\begin{center}
\scriptsize
\begin{tabularx}{\linewidth}{|X|X|X|X|X|X|X|}
\hline
\textbf{Camada} & \textbf{Responsabilidade} & \textbf{Insumo} & \textbf{Saída} &
\textbf{Dependência} & \textbf{Fronteira} & \textbf{Autoridade excluída} \\
\hline
Governança constitucional & Definir papéis, restrições, emendas,
escalonamento e condições de contestação de ordem superior & Regras
institucionais e base de emenda & Registro da restrição e da autoridade
aplicáveis & Restringe domínios locais, módulos e condições de revisão &
Restrições não delegáveis não podem ser ampliadas localmente; as regras
permanecem inspecionáveis & Nenhuma alegação de força jurídica, implantação,
infalibilidade ou execução automática \\
\hline
Domínio delimitado de governança local & Interpretar o escopo permitido para
um contexto decisório local & Restrições constitucionais e proposta local &
Decisão, configuração ou solicitação de escalonamento delimitada & Depende do
escopo constitucional e das declarações do módulo & Não pode sobrepor
restrições não delegáveis & Nenhuma autoridade soberana, ilimitada ou de
produção \\
\hline
Módulo funcional de serviço & Expressar uma capacidade institucional delimitada
e suas dependências & Escopo autorizado, tarefa, requisitos de evidência e
declarações de dependência & Proposta, resultado ou registro de falha específico
da capacidade & Restringido pela governança; inspecionado pelas camadas de
evidência e risco & Deve expor propósito, dependências, necessidades de
evidência, escalonamento e condições de falha & Não detém autoridade
constitucional, proveniência, disposição de IA ou contestabilidade \\
\hline
Evidência e proveniência & Atribuir fontes, transformações, agentes,
assistência e disposições & Artefatos de origem, declarações de autoridade,
análises, conflitos e decisões & Registro conectado de decisão e evidência &
Recebe registros de todas as camadas decisórias e apoia revisão posterior & A
conexão de registros não estabelece verdade, integridade, legitimidade ou
responsabilização & Nenhuma validação automática, julgamento ou garantia de
integridade \\
\hline
Assistência de IA & Fornecer análise consultiva delimitada e atribuível &
Tarefa declarada, material de origem e restrições & Saída de IA com
proveniência e disposição humana & Depende de insumos de evidência e revisão
humana; sujeita a limites de risco & Autoridade e responsabilidade humanas
permanecem explícitas & Nenhuma decisão soberana, execução ou julgamento
vinculante \\
\hline
Limite de risco e contestabilidade & Tornar visíveis gatilhos de revisão,
condições de contestação, necessidades de independência e disposições possíveis
& Ação proposta, limite aplicável, evidência contestada e conflitos & Registro
de revisão, recusa, escalonamento, correção, reversão ou ausência de mudança &
Depende do escopo constitucional e de evidência reconstruível & Um caminho de
contestação deve ser identificável e separável da disposição original &
Nenhuma alegação de procedimento permanente, reparação efetiva ou
responsabilização garantida \\
\hline
\end{tabularx}
\end{center}
\end{landscape}

\subsection{Camada de governança constitucional}

A camada de governança constitucional é uma superfície conceitual de
restrições que define regras de ordem superior: papéis, direitos de decisão,
proibições, procedimentos de emenda, caminhos de escalonamento e condições de
contestação. Seu propósito é tornar visíveis as restrições institucionais não
delegáveis antes do início da configuração local e enquadrar como uma revisão
posterior pode distinguir uma ação permitida de outra fora do escopo.

Nesta proposta, a camada estrutura fronteiras de apoio à responsabilização, em
vez de servir como constituição implantada, instrumento jurídico, autoridade
ativa ou órgão de execução. Uma ação revisável depende, portanto, tanto de
justificação institucional quanto de qualquer autorização técnica que uma
implementação futura possa usar; nenhuma delas, isoladamente, estabeleceria
legitimidade.

Como as próprias restrições constitucionais podem ser incompletas, ambíguas ou
excessivamente amplas, elas precisam permanecer como objetos de design
inspecionáveis, com condições delimitadas de revisão, emenda e contestação. O
modelo não trata a camada constitucional como infalível ou além de
contestação.

\subsection{Camada de domínio delimitado de governança local}

Como referência conceitual, e não como equivalência de implementação, a
governança policêntrica descreve múltiplos centros decisórios que podem ser
formalmente distintos, porém interdependentes, e operar entre camadas de
governança aninhadas \citep{ostrom2010beyond}.
A camada de domínio delimitado de governança local representa um contexto
decisório de escopo local que opera sob restrições constitucionais. Um domínio
delimitado de governança local pode configurar processos locais, selecionar
módulos funcionais de serviço e alocar tarefas decisórias apenas dentro do
escopo permitido pela camada constitucional. A documentação institucional do
Axodus pode usar o termo Tenant para esses domínios; este artigo usa a
expressão de domínio de governança para evitar implicar tenancy SaaS
implantada, entidades-clientes ativas, comunidades operantes ou infraestrutura
de produção.

Um domínio delimitado de governança local não pode ampliar ou sobrepor
restrições constitucionais não delegáveis. Quando a interpretação local de uma
ação proposta entra em conflito com uma fronteira constitucional, a
arquitetura trata essa ambiguidade como condição de revisão que exige
documentação e, quando necessário, escalonamento antes da disposição final.

\subsection{Camada de módulo funcional de serviço}

Como referência conceitual, e não como validação desta taxonomia, a literatura
sobre modularidade descreve sistemas complexos construídos a partir de
subsistemas menores que podem ser projetados separadamente e funcionar em
conjunto sob regras de design compartilhadas \citep{baldwin2000design}.
Módulos funcionais de serviço são capacidades institucionais abstratas e
delimitadas pelas quais funções selecionadas podem ser configuradas, limitadas
e governadas. O artigo não apresenta um catálogo de produtos. Em vez disso,
trata um módulo funcional de serviço como unidade conceitual que deve declarar
propósito, insumos, saídas, escopo autorizado e dependências decisórias,
requisitos de evidência, dependências, condições de escalonamento e
comportamento de falha.

Um módulo funcional de serviço é, portanto, uma superfície delimitada de
capacidade, não o detentor de autoridade constitucional, proveniência de
evidência, assistência de IA ou contestabilidade. Essas responsabilidades
permanecem em camadas transversais separadas que restringem, inspecionam e
reconstroem decisões específicas de módulos.

A documentação institucional do Axodus pode usar o termo núcleos de serviço em
contextos institucionais mais amplos; este artigo usa módulos funcionais de
serviço como abstração voltada à arquitetura para evitar implicar catálogo de
produtos, inventário de implementação ou decomposição de sistema implantado.

Uma taxonomia funcional mínima pode incluir registros de autoridade e papéis,
registros de evidência, registros de justificativa decisória, registros de
conflitos e divulgações, caminhos de revisão e contestação, controladores de
limites de risco e apoio à prontidão de participação como tipos de componentes
associados a capacidades governadas por módulos. São tipos de componentes, não
alegações sobre um inventário definitivo de módulos do Axodus ou seu estado de
implementação.

\subsection{Camada de evidência e proveniência}

A camada de evidência e proveniência conecta propostas, materiais de apoio,
declarações de autoridade, conflitos declarados, análises geradas por IA,
decisões e revisões subsequentes. Seu objetivo de design é a
reconstrutibilidade: um revisor deve poder perguntar o que era conhecido, quem
estava autorizado, qual configuração foi proposta, qual assistência foi usada
e por que um resultado se seguiu.

No nível conceitual, documentos de origem e saídas podem ser tratados como
entidades, etapas de análise e revisão como atividades, e participantes
humanos ou automatizados como agentes. Relações de derivação e atribuição
podem então conectar uma saída a seus insumos e papéis responsáveis
\citep{w3c2013provdm}. Esses conceitos especificam o que um registro revisável
deve expressar para que um revisor posterior possa reconstruir o caminho do
material de origem até a disposição. Eles não provam, por si só, que um
registro seja verdadeiro, completo, resistente a adulteração ou
responsabilizável no sentido normativo mais forte.

Assim, a camada apoia inspeção posterior, e não validação automática. Ela não
prescreve sistema de armazenamento, prova criptográfica, infraestrutura de
evidência implementada ou garantia de integridade, e não transforma a conexão
entre registros em julgamento completo sobre se a decisão foi
substantivamente justificada.

\subsection{Camada de assistência de IA}

A assistência de IA permanece consultiva. Ela pode apoiar tarefas delimitadas,
como comparação de documentos, resumo de propostas, detecção de
inconsistências, apoio à análise constitucional, verificação adversarial e
apoio à revisão orientada ao planejamento. Cada uso deve declarar tarefa,
material de origem, saída, incerteza quando disponível, revisor humano e
disposição, para que revisores posteriores possam inspecionar como a análise
consultiva entrou no registro. Humanos retêm autoridade decisória e
responsabilidade.

Esta camada é intencionalmente abstrata quanto a papéis. O artigo não usa
rótulos de agentes nomeados, não alega autoridade autônoma de execução e não
afirma que a assistência de IA execute governança, emita julgamentos
vinculantes ou substitua a revisão institucional humana.

\subsection{Camada de limite de risco e contestabilidade}

A camada de limite de risco e contestabilidade é uma superfície conceitual
para tornar visível quando revisão adicional, recusa, documentação ou um
caminho de contestação podem ser necessários. Nesse modelo, uma fronteira deve
identificar o papel relevante, o gatilho de revisão, a evidência exigida para
inspeção e a condição sob a qual uma ação proposta excede seu escopo permitido.
Contestabilidade diz respeito, portanto, a uma decisão poder ser reconstruída,
questionada e reexaminada em princípio, não a a arquitetura já operar um
mecanismo permanente de contestação.

Um registro minimamente contestável deve identificar a decisão ou o registro
contestado, a regra ou fronteira aplicável, a evidência ou justificativa em
disputa, a condição de independência esperada para o reexame e a categoria de
disposição possível, como correção, condição de escalonamento, recusa,
reversão ou ausência de mudança. Esses são campos conceituais para
reconstrutibilidade e contestabilidade, não um procedimento permanente.

Nesse enquadramento, a independência do revisor é uma preocupação de design,
e não uma alegação sobre um órgão permanente ou papel nomeado. A arquitetura
deve permitir distinguir o contexto da decisão original das condições
posteriores de contestação ou reexame, incluindo conflitos divulgados e
situações nas quais o mesmo ator não deve ser o único julgador de um registro
contestado. A supervisão humana torna-se nominal quando o revisor carece de
tempo, contexto de autoridade, acesso à evidência ou capacidade prática de
alterar um resultado \citep{tabassi2023airmf,amershi2019guidelines}. A
literatura sobre contestabilidade por design é relevante porque trata
visibilidade da revisão, intervenção e escrutínio externo como requisitos de
design para sistemas contestáveis, e não como garantias de prevenção de
resultados prejudiciais \citep{alfrink2023contestable}.

\subsection{Fronteira de apoio à Prontidão de Participação}

A Prontidão de Participação é mantida apenas como preocupação de apoio
delimitada para preparação, adaptação, revisão e contestação em contextos
consequentes selecionados. Não é contribuição coequivalente, sistema geral de
classificação social nem filtro universal de competência. Qualquer artefato de
prontidão seria um insumo contestável e revisável, não fonte independente de
legitimidade, justiça ou autoridade. Os critérios precisariam permanecer
proporcionais, sujeitos a adaptações e abertos à contestação; exclusão,
vigilância, captura interpretativa, manipulação de testes e risco de
privacidade continuam sendo razões para restringir, redesenhar ou rejeitar o
mecanismo. Sua validade é pesquisa futura, e não uma alegação avaliada neste
artigo.

\section{Status de Protótipo}

O Axodus é apresentado neste artigo como um modelo conceitual em estágio de
protótipo. O repositório público atualmente sustenta artefatos editoriais,
registros de governança e documentação de pesquisa. Ele não fornece evidência
de implementação integrada, plataforma federada em funcionamento, operação em
produção, adoção por usuários, desempenho empírico ou controles efetivos.

A ausência dessas alegações é substantiva. A consistência arquitetural pode
ser revisada antes da implementação, mas não substitui comportamento observado.
Versões futuras devem manter uma tabela de evidências por componente que
distinga estados projetados, prototipados, testados, implantados e avaliados
independentemente.

Maturidade documental, maturidade de desenvolvimento, prontidão para produção
e autoridade institucional são determinações separadas. Essa distinção serve
apenas para evitar erro de categoria: documentação não estabelece
implementação, implementação não estabelece autorização de produção e
maturidade não estabelece autoridade legítima ou capacidade operacional. O
artigo não mantém Nível-L ou Nível-D como construtos científicos independentes,
não atribui níveis a componentes e não apresenta uma taxonomia de maturidade
como instrumento validado.

A redução de escopo permanece um requisito de design. Protótipos iniciais
devem se concentrar em decisões para as quais reconstrução e contestação sejam
materialmente importantes, usar um registro delimitado de decisão e evidência,
separar saídas consultivas de IA de disposições humanas e remover elementos
que não melhorem a inspecionabilidade.

\section{Plano de Avaliação}

O programa de avaliação é organizado em torno da arquitetura integrada,
permanecendo controlado, conceitual, sintético e baseado em artefatos. Ele
avalia se a decomposição proposta torna autoridade, dependências, registros,
assistência consultiva e condições de contestação mais inspecionáveis. Não
avalia uma instituição em operação nem estabelece legitimidade, justiça,
responsabilização, efetividade da contestabilidade ou reparação.

O estágio inicial do estudo usaria cenários sintéticos de governança local
delimitada, pacotes controlados de evidências, registros simulados de
assistência de IA, erros inseridos e uma interface de revisão de alta
fidelidade em \emph{mock-up} ou \emph{wizard-of-oz}, em vez de um contexto
institucional ativo. Um protocolo pré-registrado definiria tarefas dos
revisores e critérios do estudo. Esses são desenhos propostos, não resultados
relatados.

\paragraph{RQ1: Rastreabilidade e reconstrução.} A estrutura em camadas de
decisão e evidência proposta ajuda revisores independentes a reconstruir
decisões de governança local delimitada que envolvem módulos funcionais de
serviço com maior precisão e eficiência do que uma linha de base apenas
documental? Critérios ilustrativos incluem precisão e completude da
reconstrução, visibilidade da proveniência, detecção de conflitos de
autoridade, detecção de evidência ausente, tempo para reconstruir a
justificativa e discordância entre revisores.

\paragraph{RQ2: Coerência entre camadas e condições de revisão.} Revisores
independentes conseguem usar a representação de seis camadas para identificar
autoridade ausente ou conflitante, dependências não declaradas, transições
entre camadas sem suporte e se foram especificados um caminho de contestação
acessível e uma condição de independência de revisão? Critérios ilustrativos
incluem consistência entre camadas, completude das dependências, detecção de
transições sem suporte, acessibilidade do caminho de contestação e
classificação correta das categorias possíveis de disposição. Essas medidas
testam inspecionabilidade e condições de apoio à responsabilização ou à
contestabilidade; não demonstram prestação de contas institucional,
contestação efetiva, correção, reversão ou reparação.

\paragraph{RQ3: Confiança humano--IA na revisão consultiva.} Para tarefas
delimitadas de análise de propostas, como a proveniência e a disposição humana
obrigatória afetam detecção de erros, confiança calibrada, comportamento de
anulação, qualidade da justificativa e qualidade da revisão? Registros
simulados de IA incluiriam erros inseridos, casos ambíguos e recomendações
contestadas. Critérios ilustrativos incluem detecção de erros inseridos e de
saídas sem suporte, comportamento de anulação ou aceitação por condição,
qualidade da justificativa, discordância entre revisores e correção após a
exibição de informações de proveniência.

Para este artigo, uma linha de base apenas documental usa políticas e registros
estáticos sem a arquitetura de revisão em camadas proposta, estrutura explícita
de proveniência ou enquadramento delimitado do caminho de contestação.
Resultados negativos importam tanto quanto resultados positivos: a
representação pode não melhorar a reconstrução, pode aumentar a carga de
trabalho sem melhorar a detecção de erros, pode deixar inconsistências entre
camadas sem detecção, pode produzir discordância sobre condições de
contestação ou pode não melhorar a confiança calibrada.

\paragraph{Avaliação futura em estágios.} A avaliação deve avançar da inspeção
controlada de artefatos para estudos delimitados com revisores usando cenários
sintéticos e, depois, para avaliação controlada de protótipo sob protocolo
aprovado separadamente. Qualquer estudo posterior situado institucionalmente
exigiria gates adicionais de governança, ética, direito, privacidade e
evidência. A validade da Prontidão de Participação e a responsabilização
institucional mais ampla permanecem tópicos posteriores de pesquisa.

\section{Limitações e Ameaças à Validade}

Este artigo é conceitual e ainda mais amplo do que seria ideal para um artigo
externo maduro. A arquitetura pode parecer coerente porque ainda não encontrou
restrições de implementação, conflito institucional, pressão de carga de
trabalho ou condições de campo que testariam se seus caminhos propostos de
revisão são viáveis entre decisões constitucionais e decisões em domínios
delimitados de governança local.

A base pública de evidências para o Axodus está atualmente limitada à
documentação de design. Não há resultados para avaliar viabilidade,
usabilidade, legitimidade da governança, confiança calibrada, caminhos
contestáveis ou impacto social. O artigo não relata validação de campo,
validação situada institucionalmente, uso em produção, alegações de design ou
distribuição de tokens, alegações sobre operação financeira, garantias de
segurança ou conclusão jurídica ou regulatória.

A supervisão humana pode se tornar nominal quando revisores carecem de tempo,
autoridade, acesso à evidência ou poder prático para alterar o resultado. A
assistência de IA pode ser não confiável, difícil de calibrar ou difícil de
auditar sob carga de trabalho e pressão de tempo realistas. Contestabilidade
fraca ou caminhos de contestação mal projetados podem tornar a rastreabilidade
visível sem produzir prestação de contas, capacidade corretiva, reparação
efetiva ou responsabilização institucional significativa.

Mecanismos de Prontidão de Participação apresentam riscos particulares de
exclusão, vigilância, captura curricular e falsa confiança. A separação
modular também pode ocultar acoplamento sistêmico, enquanto restrições
constitucionais podem ser excessivamente rígidas ou ambíguas para governança
local efetiva. A arquitetura proposta pode ser complexa demais para ser
governada na prática sem redução substancial de escopo.

O cenário sintético é explicativo, e não empírico. O vocabulário de
proveniência não estabelece integridade de registros ou responsabilização, e
as medidas propostas avaliam reconstrutibilidade, inspecionabilidade entre
camadas, condições de revisão e confiança delimitada, em vez de legitimidade,
justiça, contestabilidade efetiva, reparação ou responsabilização
institucional plena. Decomposições alternativas podem se mostrar mais claras
ou úteis do que o modelo de seis camadas. A avaliação situada
institucionalmente também pode expor conflitos, incentivos e efeitos de carga
de trabalho que um artefato controlado não consegue reproduzir.

O artigo não fornece aconselhamento financeiro, de segurança, de investimento
ou jurídico. Sua terminologia pode não se mapear de forma uniforme em todas as
jurisdições ou comunidades de pesquisa.

\section{Agenda de Pesquisa}

O trabalho de curto prazo deve priorizar os artefatos de derivação e interface,
cenários sintéticos que envolvam domínios delimitados de governança local e
tarefas de revisão adequadas a estudos de reconstrução, coerência entre
camadas, condições de contestação e confiança delimitada. O aprofundamento
posterior da literatura sobre qualidade da governança, independência do
revisor, caminhos empíricos de contestação institucional e validade da
Prontidão de Participação permanece requisito futuro antes que qualquer
revisão externa mais ampla seja considerada.

\section{Conclusão}

O Axodus é proposto como uma arquitetura conceitual em estágio de protótipo
para examinar como a governança assistida por humanos e IA pode se tornar mais
rastreável, reconstruível, revisável e favorável à responsabilização e à
contestabilidade em instituições digitais modulares federadas. Sua
contribuição não é um catálogo de produtos, uma alegação de estado de
implementação ou a afirmação de uma decomposição unicamente correta, mas um
arranjo inspecionável de seis camadas que conecta restrições constitucionais,
domínios delimitados de governança local, módulos funcionais de serviço,
evidência e proveniência, assistência consultiva de IA e caminhos de revisão.

O próximo passo de pesquisa é instanciar a representação conceitual em um
ambiente experimental controlado e testar se revisores independentes conseguem
reconstruir decisões, identificar inconsistências entre camadas, inspecionar
condições de contestação e avaliar a confiança delimitada em IA. Esses
resultados testariam as condições de apoio propostas pela arquitetura; não
estabeleceriam, de forma independente, legitimidade, reparação efetiva ou
responsabilização institucional.

\renewcommand{\bibpreamble}{%
\noindent As referências são mantidas em \texttt{references.bib}. Registros de
verificação de todas as entradas citadas são mantidos nos controles
bibliográficos do repositório; a rechecagem final dos metadados é necessária
antes de qualquer submissão, divulgação pública ou circulação externa.\par\medskip}
\bibliographystyle{plainnat}
\bibliography{references}

\end{document}
